%------------------------------------------------------------------------------%
\documentclass[fleqn,utf8,aspectratio=169,12pt]{beamer}

\usepackage{gaal}

\title{Sistemas Lineares Homogêneos}

%------------------------------------------------------------------------------%
\begin{document}

\addtitle

%------------------------------------------------------------------------------%
\section{Sistemas Lineares Homogêneos}

%------------------------------------------------------------------------------%
\begin{frame}{Sistemas Lineares Homogêneos}
  Um sistema linear \emph{homogêneo} tem todos os
  \emph{termos independentes iguais a zero}
  \pause
  \[
    \left\{\begin{array}{lll}
      a_{11} x_1 + \cdots + a_{1n} x_n    & \!\! = & \emph{0} \\
      a_{21} x_1 + \cdots + a_{2n} x_n    & \!\! = & \emph{0} \\
      \hphantom{a_{21} x_1 + } \quad \vdots &        &   \\
      a_{m1} x_1 + \cdots + a_{mn} x_n    & \!\! = & \emph{0}
    \end{array}\right.
  \]

  \pause
  Na forma matricial
  \[
    A_{m\times n} \; \mathbf{x}_{n\times 1} = \emph{\mnula}_{m\times 1}
  \]
\end{frame}

%------------------------------------------------------------------------------%
\begin{frame}{Existência de Solução}
  Todo sistema homogêneo possui pelo menos a \emph{solução trivial}
  \pause
  \[
    \mathbf{x}_{n\times 1}
    =
    \mnula_{n\times 1}
    \pause
    =
    \begin{bmatrix}
      0      \\
      \multicolumn{1}{c}{\vdots} \\
      0
    \end{bmatrix}_{n\times 1}
  \]

  \pause
  pois
  \[
    A\mathbf{x}
    \pause
    = A\mnula
    \pause
    = \mnula
  \]

  \emph{Atenção:} As ordens das matrizes precisam ser compatíveis
\end{frame}

%------------------------------------------------------------------------------%
\begin{frame}{Soluções de um Sistema Homogêneo}
  Um sistema homogêneo pode possuir
  \pause
  \begin{itemize}
    \item somente a solução trivial
    \pause
    \item infinitas soluções
  \end{itemize}

  \pause
  \bigskip
  Qualquer solução
  \;\(
    \mathbf{x} \neq \mnula
  \)\;
  \pause
  é chamada \emph{solução não trivial}
\end{frame}

%------------------------------------------------------------------------------%
%------------------------------------------------------------------------------%
\begin{question}
  Determine as soluções do sistema
  \[
    \systeme{
       x + 2y = 0,
      3x -  y = 0
    }
  \]

  \bigskip
  \pause
  Matriz aumentada
  \[
    \begin{amatrix}{2}
      1 & 2  & 0 \\
      3 & -1 & 0
    \end{amatrix}
  \]
\end{question}

%------------------------------------------------------------------------------%
\begin{resolution}{Escalonamento}
  \begin{columns}
  \column{0.3\linewidth}
    \[\;\;\;
    \left[\begin{array}{cc|c}
      1 & 2  & 0 \\
      3 & -1 & 0
    \end{array}\right]
    \]
    \pause
  \column{0.7\linewidth}
    Eliminamos os elementos abaixo do primeiro pivô
    \\ \pause
    \(
      L_2 \leftarrow L_2 - 3 L_1
    \)
  \end{columns}
  \vspace{-\baselineskip}
  \[
    \pause
    L'_2
    \;=\;
    L_2 - 3 L_1
    \pause
    \;=\;
    \left[\begin{array}{cc|c}
      3 & -1 & 0
    \end{array}\right]
    - 3
    \left[\begin{array}{cc|c}
      1 & 2  & 0
    \end{array}\right]
    \pause
    \;=\;
    \left[\begin{array}{cc|c}
      0 & -7 & 0
    \end{array}\right]
  \]
  \pause
  \[
    \left[\begin{array}{cc|c}
      1 & 2  & 0 \\
      0 & -7 & 0
    \end{array}\right]
  \]
\end{resolution}

%------------------------------------------------------------------------------%
\begin{resolution}{Escalonamento}
  \begin{columns}
  \column{0.3\linewidth}
    \[\;\;\;
    \left[\begin{array}{cc|c}
      1 & 2  & 0 \\
      0 & -7 & 0
    \end{array}\right]
    \]
    \pause
  \column{0.7\linewidth}
    Tornando o segundo pivô igual a 1
    \\ \pause
    \(
      L_2 \leftarrow - \frac{1}{7} L_2
    \)
  \end{columns}
  \vspace{-\baselineskip}
  \[
    \pause
    L'_2
    \;=\;
    - \frac{1}{7} L_2
    \pause
    \;=\;
    - \frac{1}{7}
    \left[\begin{array}{cc|c}
      0 & -7 & 0
    \end{array}\right]
    \pause
    \;=\;
    \left[\begin{array}{cc|c}
      0 & 1 & 0
    \end{array}\right]
  \]
  \pause
  \[
    \left[\begin{array}{cc|c}
      1 & 2 & 0 \\
      0 & 1 & 0
    \end{array}\right]
  \]
\end{resolution}

%------------------------------------------------------------------------------%
\begin{resolution}{Escalonamento}
  \begin{columns}
  \column{0.3\linewidth}
    \[\;\;\;
    \left[\begin{array}{cc|c}
      1 & 2 & 0 \\
      0 & 1 & 0
    \end{array}\right]
    \]
    \pause
  \column{0.7\linewidth}
    Eliminamos os elementos a cima do segundo pivô
    \\ \pause
    \(
      L_1 \leftarrow L_1 - 2 L_2
    \)
  \end{columns}
  \vspace{-\baselineskip}
  \[
    \pause
    L'_1
    \;=\;
    L_1 - 2 L_2
    \pause
    \;=\;
    \left[\begin{array}{cc|c}
      1 & 2 & 0
    \end{array}\right]
    - 2
    \left[\begin{array}{cc|c}
      0 & 1 & 0
    \end{array}\right]
    \pause
    \;=\;
    \left[\begin{array}{cc|c}
      1 & 0 & 0
    \end{array}\right]
  \]
  \pause
  \[
    \left[\begin{array}{cc|c}
      1 & 0 & 0 \\
      0 & 1 & 0
    \end{array}\right]
  \]

  \pause
  \bigskip
  Observe que o a última coluna da matriz aumentada sempre ficou igual a zero
\end{resolution}

%------------------------------------------------------------------------------%
\begin{resolution}{Solução}
  A única solução é
  \begin{align*}
    x & = 0 \\
    y & = 0
  \end{align*}

  \pause
  Vetorialmente
  \[
    \mathbf{x}
    \pause
    =
    \begin{bmatrix}
      x \\
      y
    \end{bmatrix}
    \pause
    =
    \begin{bmatrix}
      0 \\
      0
    \end{bmatrix}
  \]
\end{resolution}

%------------------------------------------------------------------------------%
\endinput


%------------------------------------------------------------------------------%
\begin{frame}{Forma Escalonada e Soluções}
  Se todas as incógnitas correspondem a posições de pivô
  \\ \pause
  a única solução possível é a solução trivial

  \pause
  \bigskip

  Se existir pelo menos uma variável livre
  \\ \pause
  podemos atribuir valores a essa variável e obter soluções não triviais
\end{frame}

%------------------------------------------------------------------------------%
%------------------------------------------------------------------------------%
\begin{question}
  Determine o conjunto solução do sistema
  \[
    \systeme{
       x + 2y -  z = 0,
      2x + 4y - 2z = 0,
       x +  y +  z = 0
    }
  \]

  \bigskip
  \begin{columns}[t]
  \column{0.3\linewidth}
    \pause
    Matriz aumentada
    \[
      \begin{amatrix}{3}
        1 & 2 & -1 & 0 \\
        2 & 4 & -2 & 0 \\
        1 & 1 & 1  & 0
      \end{amatrix}
    \]
  \column{0.4\linewidth}
    \pause
    Basta escalonar a matriz $A$
    \[
      A =
      \begin{bmatrix}
        1 & 2 & -1  \\
        2 & 4 & -2  \\
        1 & 1 &  1
      \end{bmatrix}
    \]
  \end{columns}
\end{question}

%------------------------------------------------------------------------------%
\begin{resolution}{Escalonamento}
  \begin{columns}
  \column{0.3\linewidth}
    \[\;\;\;
      \begin{bmatrix}
        1 & 2 & -1 \\
        2 & 4 & -2 \\
        1 & 1 &  1
      \end{bmatrix}
    \]
    \pause
  \column{0.7\linewidth}
    Eliminamos os elementos abaixo do primeiro pivô
    \\ \pause
    \(
      L_2 \leftarrow L_2 - 2 L_1
    \)
    \\ \pause
    \(
      L_3 \leftarrow L_3 - L_1
    \)
  \end{columns}
  \vspace{-\baselineskip}
  \[
    \pause
    L'_2
    \;=\;
    L_2 - 2 L_1
    \pause
    \;=\;
    \begin{bmatrix}
      2 & 4 & -2
    \end{bmatrix}
    - 2
    \begin{bmatrix}
      1 & 2 & -1
    \end{bmatrix}
    \pause
    \;=\;
    \begin{bmatrix}
      0 & 0  & 0
    \end{bmatrix}
  \]
  \[
    \pause
    L'_3
    \;=\;
    L_3 - L_1
    \pause
    \;=\;
    \begin{bmatrix}
      1 & 1 &  1
    \end{bmatrix}
    -
    \begin{bmatrix}
      1 & 2 & -1
    \end{bmatrix}
    \pause
    \;=\;
    \begin{bmatrix}
      0 & -1 & 2
    \end{bmatrix}
  \]
  \pause
  \[
    \begin{bmatrix}
      1 & 2  & -1 \\
      0 & 0  & 0  \\
      0 & -1 & 2
    \end{bmatrix}
  \]
\end{resolution}

%------------------------------------------------------------------------------%
\begin{resolution}{Escalonamento}
  \begin{columns}
  \column{0.3\linewidth}
    \[\;\;\;
      \begin{bmatrix}
        1 & 2  & -1 \\
        0 & 0  & 0  \\
        0 & -1 & 2
      \end{bmatrix}
    \]
    \pause
  \column{0.7\linewidth}
    Trocamos as linhas $L_2$ e $L_3$
  \end{columns}

  \pause
  \[
    \begin{bmatrix}
      1 & 2  & -1 \\
      0 & -1 & 2  \\
      0 & 0  & 0
    \end{bmatrix}
  \]
\end{resolution}

%------------------------------------------------------------------------------%
\begin{resolution}{Escalonamento}
  \begin{columns}
  \column{0.3\linewidth}
    \[\;\;\;
      \begin{bmatrix}
        1 & 2  & -1 \\
        0 & -1 & 2  \\
        0 & 0  & 0
      \end{bmatrix}
    \]
    \pause
  \column{0.7\linewidth}
    Tornando o segundo pivô igual a 1
    \\ \pause
    \(
      L_2 \leftarrow - L_2
    \)
  \end{columns}
  \vspace{-\baselineskip}
  \[
    \pause
    L'_2
    \;=\;
    - L_2
    \pause
    \;=\;
    -
    \begin{bmatrix}
      0 & -1 & 2
    \end{bmatrix}
    \pause
    \;=\;
    \begin{bmatrix}
      0 & 1 & -2
    \end{bmatrix}
  \]
  \pause
  \[
    \begin{bmatrix}
      1 & 2 & -1 \\
      0 & 1 & -2 \\
      0 & 0 & 0
    \end{bmatrix}
  \]
\end{resolution}

%------------------------------------------------------------------------------%
\begin{resolution}{Escalonamento}
  \begin{columns}
  \column{0.3\linewidth}
    \[\;\;\;
      \begin{bmatrix}
        1 & 2 & -1 \\
        0 & 1 & -2 \\
        0 & 0 & 0
      \end{bmatrix}
    \]
    \pause
  \column{0.7\linewidth}
    Eliminamos os elementos a cima do segundo pivô
    \\ \pause
    \(
      L_1 \leftarrow L_1 - 2 L_2
    \)
  \end{columns}
  \vspace{-\baselineskip}
  \[
    \pause
    L'_1
    \;=\;
    L_1 - 2 L_2
    \pause
    \;=\;
    \begin{bmatrix}
      1 & 2 & -1
    \end{bmatrix}
    - 2
    \begin{bmatrix}
      0 & 1 & -2
    \end{bmatrix}
    \pause
    \;=\;
    \begin{bmatrix}
      1 & 0 & 3
    \end{bmatrix}
  \]
  \pause
  \[
    \begin{bmatrix}
      1 & 0 & 3  \\
      0 & 1 & -2 \\
      0 & 0 & 0
    \end{bmatrix}
  \]
\end{resolution}

%------------------------------------------------------------------------------%
\begin{resolution}{Solução}
  Sistema reduzido
  \[
    \systeme{
      x   + 3z = 0,
        y - 2z = 0
    }
  \]

  \pause
  \bigskip
  A variável $z$ é livre
  \[
    z = t
    \qquad
    \forall\, t \in \R
  \]
\end{resolution}

%------------------------------------------------------------------------------%
\begin{resolution}{Solução}
  Solução
  \begin{align*}
    x & = -3t \\
    y & =  2t
  \end{align*}

  \pause
  Vetorialmente
  \[
    \mathbf{x}
    \pause
    \;=\;
    \begin{bmatrix}
      -3t \\
      2t  \\
      t
    \end{bmatrix}
    \pause
    \;=\;
    t
    \begin{bmatrix}
      -3 \\
      2  \\
      1
    \end{bmatrix}
  \]
\end{resolution}

%------------------------------------------------------------------------------%
\begin{resolution}{Conjunto Solução}
  O \emph{conjunto solução} é
  \[
    S
    =
    \left\{
    t
    \begin{bmatrix}
      -3 \\
      2  \\
      1
    \end{bmatrix}
    \st
    t\in\mathbb{R}
    \right\}
  \]

  \pause
  O conjunto solução forma uma reta que passa pela origem
\end{resolution}

%------------------------------------------------------------------------------%
\endinput

%------------------------------------------------------------------------------%
\begin{question}
  Determine o conjunto solução do sistema
  \[
    \systeme[sort={x, y, z, w}]{
       x +  y +  z +  w = 0,
      2x + 2y + 2z + 2w = 0
    }
  \]

  \bigskip
  \[
    A =
    \begin{bmatrix}
      1 & 1 & 1 & 1 \\
      2 & 2 & 2 & 2
    \end{bmatrix}
  \]
\end{question}

%------------------------------------------------------------------------------%
\begin{resolution}{Escalonamento}
  \begin{columns}
  \column{0.3\linewidth}
    \[\;\;\;
      \begin{bmatrix}
        1 & 1 & 1 & 1 \\
        2 & 2 & 2 & 2
      \end{bmatrix}
    \]
    \pause
  \column{0.7\linewidth}
    Eliminamos os elementos abaixo do primeiro pivô
    \\ \pause
    \(
      L_2 \leftarrow L_2 - 2 L_1
    \)
  \end{columns}
  \vspace{-\baselineskip}
  \[
    \pause
    L'_2
    \;=\;
    L_2 - 2 L_1
    \pause
    \;=\;
    \begin{bmatrix}
      2 & 2 & 2 & 2
    \end{bmatrix}
    - 2
    \begin{bmatrix}
      1 & 1 & 1 & 1
    \end{bmatrix}
    \pause
    \;=\;
    \begin{bmatrix}
      0 & 0 & 0 & 0
    \end{bmatrix}
  \]
  \pause
  \[
    \begin{bmatrix}
      1 & 1 & 1 & 1 \\
      0 & 0 & 0 & 0
    \end{bmatrix}
  \]
\end{resolution}

%------------------------------------------------------------------------------%
\begin{resolution}{Solução}
  \onslide<+->{Sistema reduzido}
  \begin{gather*}
    \onslide<+->{x + y + z + w = 0 \\}
    \onslide<+->{x = -y -z -w}
  \end{gather*}
  \onslide<+->{As variáveis $y$, $z$ e $w$ são livres}
  \begin{align*}
    \onslide<+->{y & = r & \forall\, r \in \R \\}
    \onslide<+->{z & = s & \forall\, s \in \R \\}
    \onslide<+->{w & = t & \forall\, t \in \R \\}
    \onslide<+->{x & = -r -s -t}
  \end{align*}
\end{resolution}

%------------------------------------------------------------------------------%
\begin{resolution}{Solução}
  Vetorialmente
  \[
    \mathbf{x}
    \pause
    \;=\;
    \begin{pmatrix}
      -r-s-t \\
      r      \\
      s      \\
      t
    \end{pmatrix}
    \pause
    \;=\;
    r \begin{pmatrix} -1 \\ 1 \\ 0 \\ 0 \end{pmatrix} +
    s \begin{pmatrix} -1 \\ 0 \\ 1 \\ 0 \end{pmatrix} +
    t \begin{pmatrix} -1 \\ 0 \\ 0 \\ 1 \end{pmatrix}
  \]
  \pause
  com
  \(
    r, s, t \in \R
  \)
\end{resolution}

%------------------------------------------------------------------------------%
\begin{resolution}{Conjunto Solução}
  \[
    S = \left\{
      r\, \mathbf{V}_1 +
      s\, \mathbf{V}_2 +
      t\, \mathbf{V}_3
      \st
      (r, s, t) \in \R^3
    \right\}
  \]
  com
  \[
    \mathbf{V}_1 = \begin{pmatrix} -1 \\ 1 \\ 0 \\ 0 \end{pmatrix} \qquad
    \mathbf{V}_2 = \begin{pmatrix} -1 \\ 0 \\ 1 \\ 0 \end{pmatrix} \qquad
    \mathbf{V}_3 = \begin{pmatrix} -1 \\ 0 \\ 0 \\ 1 \end{pmatrix}
  \]
\end{resolution}

%------------------------------------------------------------------------------%
\endinput


%------------------------------------------------------------------------------%
%\section{Lista Mínima}
%
%%------------------------------------------------------------------------------%
%\begin{frame}{Lista Mínima}
%  \begin{enumerate}
%    \item
%  \end{enumerate}
%
%  \vfill
%  Atenção: A prova é baseada no livro, não nas apresentações
%\end{frame}

%------------------------------------------------------------------------------%
\end{document}
%------------------------------------------------------------------------------%




